\uuid{plpb}
\titre{Matrice d'une application linéaire/Théorème du rang}
\auteur{hanani ; mzari}
\organisation{exo7}
\licence{CC BY-NC-SA 4.0}
\datecreate{2020-01-20}
\niveau{L1}
\module{Algèbre}
\chapitre{Matrice}
\sousChapitre{Matrice et application linéaire}
\difficulte{3}
\theme{Matrice d'une application linéaire, Théorème du rang}
\exo7qcmid{lille-2.matrice-application.022}

\contenu{
\begin{question}
\qclassification{108.03}
\qtags{motcle=Matrice d'une application linéaire/Théorème du rang}
Soit $A$ une matrice à coefficients réels, à $3$ lignes et $4$ colonnes. Quelles sont les assertions vraies ?
\begin{answers}
\bad{$A$ est la matrice d'une application linéaire de $\Rr^3$ dans $\Rr^4$ dans des bases de $\Rr^3$ et $\Rr^4$}.
\good{$A$ est la matrice d'une application linéaire de $\Rr^4$ dans $\Rr^3$ dans  des bases de $\Rr^4$ et $\Rr^3$}.
\bad{$A$ est la matrice d'une application linéaire de noyau  nul}.
\bad{$A$ est la matrice d'une application linéaire bijective}.
\end{answers}
\begin{explanations} $A$ est la matrice d'une application linéaire $f: \Rr^4 \to \Rr^3$ dans des bases de $\Rr^4$ et
$\Rr^3$. D'après le théorème du rang, le noyau d'une telle application est non nul.
\vskip0mm
Comme $A$ n'est pas une matrice carrée, $A$ n'est pas inversible et donc si $f$ est une application linéaire de matrice
$A$, $f$ n'est pas bijective.
\end{explanations}
\end{question}
}
